\section{文献综述}

\subsection{研究背景}

偏微分方程模型广泛出现于热传导、渗流、弹性力学、流体力学、电磁散射以及多物理场耦合问题中。在这些应用里，一个极其常见的任务并不是只求解一次方程，而是对同一类模型在不同参数下进行大量重复求解。例如，在参数辨识中需要不断试探候选参数；在最优控制与优化设计中需要反复评估目标函数；在不确定性量化中需要对参数空间进行大量采样；在实时预测与快速决策中又要求每次求解尽可能快。对于这类任务，单次求解的代价固然重要，但更关键的是总体重复求解成本，因此通常被称为 \textbf{many-query 场景}。

以参数化椭圆型方程为例，设
\[
-\nabla\cdot(a(x;\mu)\nabla u(x;\mu)) = f(x),\qquad x\in\Omega,\ \mu\in\mathcal P\subset\mathbb R^p .
\]
经过有限元离散后，可得到大规模线性系统
\[
A(\mu)u_h(\mu)=f(\mu),\qquad A(\mu)\in\mathbb R^{N\times N},\ u_h(\mu)\in\mathbb R^N,\ N\gg 1.
\]
当参数 $\mu$ 改变时，需要反复求解不同的高维线性系统。如果直接对每一个参数点都做一次高精度有限元求解，那么在 many-query 场景下，离散维数 $N$ 与参数样本数 $m$ 的叠加会迅速推高总成本。因此，如何在保证精度的前提下显著降低多参数重复求解开销，成为参数化偏微分方程数值计算中的核心问题 。

正是在这一背景下，\textbf{Reduced Basis（RB）方法} 成为参数化 PDE 的核心降阶技术之一。RB 的基本思想是：大量高维解虽然存在于 $\mathbb R^N$ 中，但对于许多参数化问题，它们往往集中在一个低维流形附近，因此可以用一个维数远小于 $N$ 的低维子空间来逼近 \cite{hesthaven2015,patera_rozza}。若记降阶基矩阵为
\[
V_r\in\mathbb R^{N\times r},\qquad r\ll N,
\]
则可用
\[
u_h(\mu)\approx V_r a(\mu)
\]
来近似高维解，再通过 Galerkin 投影得到降阶系统
\[
V_r^T A(\mu)V_r a(\mu)=V_r^T f(\mu).
\]
这样一来，昂贵的高维求解被转化为低维系统求解。RB 方法最重要的思想之一，就是通过 Offline/Online 分解将主要计算负担集中到离线阶段，而把快速响应留给在线阶段 \cite{hesthaven2015,patera_rozza}。

那么，为什么 RB 往往又要与 SVD 结合？原因在于：RB 的关键并不只是“有一个低维空间”，而是要“构造一个尽可能优的低维空间”。最经典的做法是先收集若干高精度快照
\[
U=[u_h(\mu_1),u_h(\mu_2),\dots,u_h(\mu_m)]\in\mathbb R^{N\times m},
\]
再对快照矩阵做奇异值分解
\[
U=\Phi\Sigma\Psi^T.
\]
取前 $r$ 个左奇异向量便得到 POD 基
\[
V_r=\Phi(:,1:r).
\]
这一步实际上是在 Frobenius 范数意义下寻找对所有快照的最优低秩逼近，因此 POD/SVD 在 RB 中并不是可有可无的附属步骤，而是构造高质量降阶基的核心工具。

进一步的问题是：为什么还要引入 \textbf{增量 SVD}？这是因为在实际的 RB 离线阶段，快照往往并不是一开始就全部拿到的，而是随着参数采样逐步生成的。一方面，贪婪采样、自适应采样等策略都天然带有“迭代新增样本”的结构；另一方面，在大规模仿真中，一次性存储全部快照矩阵本身就可能带来显著的内存压力。因此，若每新加入一个快照都重新对全部快照矩阵做一次 batch SVD，那么计算与存储代价都会很高 。

因此，从 many-query 的整体计算链条看，逻辑是清楚的：many-query 问题需要一个高效的参数化求解框架，因此选择 RB；RB 需要构造尽可能优的低维基，因此自然结合 POD/SVD；而 POD/SVD 在实际离线构基中又面临快照逐步增长、重复分解代价高的问题，因此进一步需要增量 SVD \cite{brand2002,brand2006,zhang2022}。

\subsection{国内外研究进展}

\subsubsection{参数化 PDE 的 Reduced Basis 与 Offline/Online 框架}

参数化 PDE 的模型降阶研究在过去二十余年发展迅速，其中最具代表性的两条主线是 POD 与 RB。Patera、Rozza 以及 Hesthaven 等人的工作系统奠定了参数化方程降阶求解的理论与算法框架：一方面，解流形在适当条件下具有较低的 Kolmogorov 宽度，说明低维近似具有理论可行性；另一方面，Offline/Online 分解将计算量大的部分集中在离线阶段，将快速求解保留到在线阶段，使 many-query 计算成为可能 \cite{patera_rozza,hesthaven2015}。

具体而言，RB 方法的优势在于它并不追求在整个高维空间内逼近，而是针对参数族解的实际分布构造一个问题相关的近似空间。对于许多椭圆型和抛物型问题，参数变化虽然会改变系数矩阵与解向量，但其主导模态常常只占据一个较低维的子空间，因此 RB 可以用很小的 $r$ 达到较高精度。与此同时，若系统具有仿射参数分解，还可以进一步预组装在线阶段的矩阵和向量，从而将在线复杂度降到与 $r$ 而非 $N$ 相关 \cite{hesthaven2015,patera_rozza}。

不过，RB 的瓶颈也非常明确：离线构基。构基质量直接决定在线精度，而构基代价又往往主导整体成本。特别是在样本数较大、单次高精度求解昂贵、快照维数很高时，如何高效地从快照中提取主导信息就成为核心问题。这也是后续 POD/SVD 与增量 SVD 被引入 RB 框架的根本原因 \cite{hesthaven2015,fareed2018}。

\subsubsection{POD/SVD 及其在降阶中的作用}

Proper Orthogonal Decomposition 本质上是一个最优低秩逼近方法。对于快照矩阵
\[
U=[u_1,\dots,u_m]\in\mathbb R^{N\times m},
\]
其截断 SVD
\[
U\approx U_r=\Phi_r\Sigma_r\Psi_r^T
\]
满足
\[
\|U-U_r\|_F^2=\sum_{i>r}\sigma_i^2,
\]
即在所有秩不超过 $r$ 的矩阵中，截断 SVD 给出了最小的 Frobenius 误差。这意味着，如果用前 $r$ 个左奇异向量张成的空间去逼近快照集合，那么从总体均方误差意义看，它是最优的 \cite{golub2013matrix}。

在 PDE 降阶中，POD 有两层意义。第一层是数据压缩：把海量高维快照压缩到少数主模态上。第二层是物理解释：奇异值衰减速度反映了解流形的有效维数，主导模态往往对应主要物理结构。对于 many-query 问题，如果快照在主模态方向上高度集中，那么说明用低维基来代替全维求解具有很强可行性 \cite{hesthaven2015,golub2013matrix}。

然而，经典 POD/SVD 也有天然缺陷。它是典型的 batch 处理方式：需要在构造好全部快照矩阵之后一次性做分解。在离线阶段，这种模式对内存与算力都不友好。特别是在有限元背景下，$N$ 很大，而参数采样数 $m$ 也可能逐步增多，此时反复对大矩阵做完整 SVD 的代价十分显著。

\subsubsection{面向 PDE 数据的加权内积与 core SVD}

对于一般的数据矩阵，标准 Euclidean 内积下的 SVD 已足够。但在有限元离散的 PDE 数据里，更自然的内积通常并不是标准内积，而是由质量矩阵或刚度矩阵诱导的加权内积。例如在有限元空间中，若 $M$ 为质量矩阵，则常考虑
\[
\langle u,v\rangle_W=u^T W v,\qquad W=M.
\]
这意味着正交性、范数与能量贡献的定义都应当相对于 $W$ 而不是单位矩阵。若忽略这一点，所构造的 POD 基就可能与连续问题中的自然能量结构不一致。

为了解决这一问题，Fareed、Singler、Zhang 与 Shen 等人发展了适用于 PDE 仿真数据的 \textbf{incremental POD / core SVD} 框架。在这一框架下，不再简单追求 $Q^TQ=I$，而是要求
\[
Q^T W Q=I,\qquad R^T R=I,
\]
并将
\[
U=Q\Sigma R^T
\]
称为 core SVD。Zhang (2022) 也在文章中沿用了这一定义，并明确指出：若数据来自 Galerkin 型 PDE 仿真，则直接应用标准 Brand 算法会牵涉到 Cholesky 分解，而通过 core SVD 可以更自然地在 weighted inner product 下工作 \cite{fareed2018,zhang2022}。

这一步对基于有限元快照的降阶问题尤其关键，因为研究对象不是一般图像流或机器学习数据，而是有限元离散下的参数化 PDE 快照，因此加权内积框架下的 core SVD 比标准 Euclidean SVD 更符合问题本身的数值结构。

\subsubsection{Brand 增量 SVD 算法及其意义}

增量奇异值分解的核心思想，是在已有矩阵奇异值分解结果的基础上，随着新列向量不断加入，快速更新当前分解，而不必每次都对扩展后的整体矩阵重新执行一次 batch SVD。对于快照逐步生成的 POD/RB 离线阶段，这种“边到达、边更新”的机制尤其重要，因此 Brand 的工作通常被视为增量 SVD 研究中的经典起点 \cite{brand2002,brand2006,zhang2022}。

在有限元离散所诱导的加权内积框架下，设
\[
\langle u,v\rangle_W=u^T W v,
\]
其中 $W$ 通常为质量矩阵或其他对称正定权矩阵。若矩阵 $U_\ell\in\mathbb R^{m\times \ell}$ 的前 $\ell$ 列已经得到了一个截断的 core SVD
\[
U_\ell = Q\Sigma R^T,\qquad Q^T W Q = I,\quad R^T R = I,
\]
其中 $Q\in\mathbb R^{m\times k}$，$\Sigma\in\mathbb R^{k\times k}$，$R\in\mathbb R^{\ell\times k}$，且 $k\le \ell$，那么 Brand 型算法的目标就是在加入新列 $u_{\ell+1}$ 后，快速构造
\[
U_{\ell+1}=[U_\ell,\ u_{\ell+1}]
\]
的更新分解 \cite{brand2002,zhang2022}。

这一算法可以分为“初始化—更新—重正交化”三个层次。

\paragraph{初始化}

若只考虑首个非零列向量 $u_1$，则初始化非常直接。此时
\[
\Sigma = (u_1^T W u_1)^{1/2},\qquad
Q=u_1\Sigma^{-1},\qquad
R=1.
\]
于是便得到单列矩阵 $U_1$ 的初始 core SVD。这个初始化步骤的本质，是先用第一个快照建立一个一维的加权正交基方向，并把对应的尺度信息吸收进奇异值 $\Sigma$ 中 \cite{zhang2022}。

\paragraph{更新}

假设当前已有
\[
U_\ell = Q\Sigma R^T.
\]
当新快照 $u_{\ell+1}$ 到来时，首先将其投影到当前左奇异子空间 $\mathrm{span}(Q)$ 上，并计算残差：
\[
d = Q^T(Wu_{\ell+1}),\qquad
e = u_{\ell+1}-Qd,\qquad
p=(e^TWe)^{1/2}.
\]
这里，$d$ 表示新列在当前基上的坐标，$e$ 表示新信息中无法被现有子空间表示的部分，$p$ 则衡量该残差的加权范数。若 $p>0$，进一步令
\[
\tilde e = e/p.
\]
这样便可写出如下恒等分解：
\[
U_{\ell+1}
=
[Q,\ \tilde e]
\begin{bmatrix}
\Sigma & Q^TWu_{\ell+1}\\
0 & p
\end{bmatrix}
\begin{bmatrix}
R & 0\\
0 & 1
\end{bmatrix}^{T}.
\]
记中间小矩阵为
\[
Y=
\begin{bmatrix}
\Sigma & d\\
0 & p
\end{bmatrix},
\]
则更新问题便被转化为求一个只有 $(k+1)\times(k+1)$ 规模的小矩阵 $Y$ 的 SVD。设
\[
Y = Q_Y\Sigma_Y R_Y^T,
\]
则理论上就可以更新得到
\[
Q^+=[Q,\ \tilde e]Q_Y,\qquad
\Sigma^+=\Sigma_Y,\qquad
R^+=
\begin{bmatrix}
R & 0\\
0 & 1
\end{bmatrix}R_Y.
\]
这正是 Brand 型增量更新的核心：把对大矩阵 $U_{\ell+1}$ 的重新分解，等价转化成对一个低维边界矩阵 $Y$ 的分解，因此显著降低了更新代价 \cite{brand2002,zhang2022}。

\paragraph{截断策略}

在实际应用中，增量 SVD 往往并不追求完整秩，而更关注主导模态，因此需要配合截断机制。Brand 型算法中至少存在两类自然截断：

第一类是\textbf{残差截断}。如果
\[
p<\mathrm{tol},
\]
说明新快照 $u_{\ell+1}$ 与已有子空间非常接近，此时可近似认为 $p=0$、$\tilde e=0$，即新增样本并未真正带来新的独立方向。于是只需更新已有模态，不必增加秩。此时常写成
\[
Q \leftarrow QQ_Y(1\!:\!k,1\!:\!k),\qquad
\Sigma \leftarrow \Sigma_Y(1\!:\!k,1\!:\!k),
\]
\[
R \leftarrow
\begin{bmatrix}
R & 0\\
0 & 1
\end{bmatrix}
R_Y(:,1\!:\!k).
\]

第二类是\textbf{奇异值截断}。若更新后最后若干个奇异值小于给定阈值，则可将其丢弃，仅保留前 $r$ 个主导模态，以控制秩增长并抑制机器精度附近的伪小奇异值。这一策略在面向 PDE 快照的实际 POD/RB 构基中尤其重要，因为目标往往是稳定提取低维主子空间，而不是恢复所有数值上微弱的方向 \cite{fareed2018,fareed2020,zhang2022}。

\paragraph{整体算法结构}

综合来看，Brand 型增量 SVD 的“完全增量版”可以概括为如下流程：先用首列初始化当前分解；随后对每一个新到达的列向量执行更新；每一步更新后，再根据需要进行重正交化修正。其结构可以形式化地写成：
先取 $u_1$ 建立初始的 $(Q,\Sigma,R)$；然后对 $\ell=2,\dots,n$，依次读入 $u_\ell$，执行更新，再执行重正交化，最终返回更新后的三元组 $(Q,\Sigma,R)$ \cite{brand2002,brand2006,zhang2022}。

从 RB/POD 离线构基的角度看，这一算法的重要意义在于：它使得快照矩阵可以按列流式到达，而低维基空间可以同步更新，从而避免“每来一个快照就重做一次 batch SVD”的高额代价。因此，在 many-query 场景下，它不仅是一个数值线性代数技巧，更是连接“高精度快照生成”和“低维缩减基提取”的关键中间环节。

\subsubsection{Brand 2006 及后续改进：从直接更新到间接更新}

尽管上述直接更新形式已经大幅优于反复执行整体 SVD，但如果每一步都显式构造
\[
Q^+=[Q,\tilde e]Q_Y,\qquad
R^+=
\begin{bmatrix}
R & 0\\
0 & 1
\end{bmatrix}R_Y,
\]
那么随着样本数不断增长，左右奇异向量矩阵仍会越来越大，更新成本也会逐步上升。Brand 在后续工作中进一步关注低秩修改问题，试图减少对外层大矩阵的频繁操作，把更多更新压缩到低维空间内部完成 \cite{brand2006}。

沿着这一思路，后续研究逐渐发展出一种更细致的分层结构：将奇异向量更新拆成“大外层矩阵 + 小正交因子”的乘积形式。这样，大矩阵主要负责记录当前子空间的外壳，而频繁发生的旋转与修正则尽可能留在小矩阵内部完成。Zhang (2022) 在此基础上进一步分析了这种结构在 weighted inner product 下的数值行为，并指出：从复杂度角度看，把更新尽量压缩到小矩阵层面是合理的；但从稳定性角度看，必须仔细区分“哪些正交因子容易失去正交性，哪些则相对稳定”，否则算法虽然快，却未必可靠 \cite{zhang2022}。

因此，从 Brand 2002 到 Brand 2006，再到 Zhang 2022，可以看到增量 SVD 的研究主线已经不再只是“如何更新”，而逐渐转向“如何在高效更新的同时维持数值稳定与正交结构”。这也正是该方法从一般数据流处理走向 PDE 降阶构基时必须面对的问题。

\paragraph{增量 POD 的优势与思考}

与传统 batch POD 相比，增量 POD 的首要优势在于其\textbf{更符合快照逐步生成的离线构基过程}。在参数化偏微分方程的 many-query 场景中，高精度快照往往不是一次性全部获得，而是随着参数采样、贪婪选点或自适应策略逐步产生。此时，若每增加一个快照都重新对全部快照矩阵执行一次 batch SVD，则不仅计算代价高，而且要求持续保存全部历史快照，带来较大的存储负担。增量 POD 则可以在已有分解基础上对新到达的数据进行递推更新，使低维基空间随样本增长同步演化，因此在算法结构上更适合实际的离线阶段。

第二，增量 POD 在\textbf{内存占用与可扩展性}方面具有明显优势。对于有限元离散得到的高维快照，状态维数 $N$ 往往很大，而样本数 $m$ 也可能随着参数探索不断增加。传统 batch POD 通常需要保留完整快照矩阵 $U\in\mathbb R^{N\times m}$，当 $N$ 和 $m$ 同时较大时，存储压力会迅速增大。相比之下，增量 POD 只需维护当前的低秩分解结果及少量新增信息，而不必始终回看全部历史数据，因此更适合大规模、流式或长期累积的数据环境。

第三，增量 POD 更有利于\textbf{与参数采样策略和在线构基过程耦合}。在实际 reduced basis 方法中，快照的选取通常并非固定不变，而是与误差估计、贪婪算法、自适应加点等策略紧密相关。增量 POD 的递推更新机制使得“新增一个参数样本—更新一次缩减基”这一过程变得自然可行，从而更容易嵌入完整的 RB 离线流程。这说明增量 POD 的意义并不只是“把 SVD 算得更快”，而是在 many-query 框架下提供了一种更加灵活的基空间构造方式。

当然，增量 POD 的优势并不是无条件的。与 batch POD 相比，它最大的代价在于\textbf{数值稳定性问题更加突出}。由于更新过程是逐步递推的，舍入误差会在多次迭代后不断累积，进而导致正交性退化、奇异值扰动以及截断误差传播等问题；在有限元加权内积框架下，这一问题又会进一步加剧。因此，增量 POD 的真正关键不只是“增量更新”本身，而在于如何在更新效率与数值稳定性之间取得平衡。换言之，增量 POD 的研究重点应当从“能否递推更新”进一步转向“如何稳定地递推更新”。


\subsection{存在问题}

\subsubsection{batch SVD 的重复分解与存储开销}

首先，最直接的问题来自 batch SVD 本身。对于一个 $N\times m$ 的快照矩阵，一次完整分解就已不便宜；若在离线阶段每新增一个快照都重做一次分解，总成本会随样本数迅速增长。更重要的是，batch SVD 通常要求同时保存全部快照数据，这在高维有限元场景下意味着明显的内存压力。因此，对于快照逐步生成的 RB 离线阶段而言，batch SVD 在计算模式上并不理想 \cite{brand2002,brand2006,fareed2018}。

\subsubsection{Brand 增量 SVD 的正交性退化问题}

相比 batch SVD，Brand 型增量 SVD 的最大优势在于每一步更新代价低；但它同时也带来了最核心的数值隐患，即\textbf{正交性退化问题}。理论上，如果在精确算术下每一步都满足
\[
Q^T W Q=I,\qquad R^TR=I,
\]
则更新后的分解始终保持正交结构。然而在实际浮点运算中，矩阵乘法、归一化和重复迭代会不断引入舍入误差；当更新次数达到上千甚至上百万次时，这些误差会逐渐积累，使本应正交的向量组偏离正交条件 \cite{brand2006,zhang2022}。

这一问题之所以严重，是因为它会直接破坏增量 SVD 的数学含义。首先，一旦
\[
Q^T W Q\neq I,
\]
当前分解就不再是严格意义上的 core SVD，相应奇异值也不再准确反映真实谱结构。这样一来，以奇异值大小为依据的截断策略可能失效。其次，POD/RB 的核心是利用主导模态张成一个高质量低维子空间；若基向量之间已经失去正交性，则不同模态可能发生污染与冗余，从而降低缩减基质量，进一步影响在线阶段的精度与稳定性。最后，在有限元加权内积下，这一问题往往比标准 Euclidean 情形更棘手，因为重正交化必须针对 $W$-内积执行，代价显著增加 \cite{fareed2018,zhang2022}。

Brand 型算法的完整流程中通常都会显式加入一个\textbf{重正交化步骤}。以 weighted Gram--Schmidt 为例，其基本思想是：当检测到正交性误差超过阈值时，依次对当前基向量做加权正交投影与归一化修正，即对第 $i$ 个向量减去其在前面各向量方向上的 $W$-投影，再按 $W$-范数归一化。这个过程理论上可以恢复当前基的加权正交性，但在实践中代价很高，尤其当基维数和更新次数都较大时，重正交步骤可能占据算法的大部分耗时 \cite{zhang2022}。

Zhang (2022) 的数值结果清楚地说明了这一点：在标准内积情形 $W=I$ 下，算法的整体代价尚可接受；但在有限元质量矩阵情形 $W=M$ 下，重正交化几乎成为主要耗时来源。文中甚至专门用
\[
E_W=\|I-Q^TWQ\|
\]
来衡量加权正交误差，并展示了在长时间迭代后，即使实施了重正交，正交性仍可能持续退化 \cite{zhang2022}。这说明问题并不只是“要不要重正交”，而是 Brand 型更新结构本身就会持续制造需要修补的正交性误差。

\subsubsection{Brand 开放问题与 Zhang (2022) 的回答}

Brand 在 2006 年提出了一个非常关键的开放问题：为了保证总体数值精度，到底需要多频繁地执行重正交化？这一问题的难点不在于“理论上可不可以重正交”，而在于：如果一味地频繁重正交，算法效率优势会被大幅削弱；但如果重正交不及时，奇异向量的正交结构又会逐渐崩塌 \cite{brand2006}。

Zhang (2022) 对这一开放问题的回答，并不是简单给出一个固定的“每多少步重正交一次”的经验法则，而是更深入地分析了 Brand 型分层更新结构中不同部分的数值脆弱性。其核心观点是：并不是所有正交因子都会同等程度地失去正交性。小规模内部因子由于维度低、连乘次数有限，往往可以在合理精度下保持稳定；真正容易快速退化的，反而是不断扩展的外层大矩阵。因此，算法设计的重点不应放在对所有小矩阵做无差别重正交，而应集中处理最外层、最脆弱的大尺度正交结构 \cite{zhang2022}。

更具体地说，Zhang 提出应尽量避免反复形成并连乘大量小正交矩阵，而把必要的正交修正集中到新增方向相对于当前大基矩阵的递归加权正交化上。这样做的结果是：一方面保留了增量更新“小矩阵计算为主”的效率优势；另一方面又减少了由长链式正交矩阵乘法带来的正交性侵蚀。这实际上把 Brand 原先的开放问题从“多久重正交一次”转化为“应当对哪一层进行重正交”，从而给出了更具结构性的答案 \cite{zhang2022}。

对于基于有限元快照的 PDE 模型降阶而言，这一点尤其重要。因为在加权内积框架下，重正交不仅昂贵，而且直接关系到 POD/RB 基的物理一致性与数值稳定性。因此，如何把 Brand 型增量更新与更稳健的重正交策略结合起来，正是当前这类课题最自然、也最有价值的研究切入点之一。

\subsection{研究展望}

结合上述研究进展与存在问题，可以看到“基于增量奇异值分解的偏微分方程模型降阶方法”具有清晰而具体的研究空间。未来可从以下几个方向深入展开。

第一，\textbf{在参数化椭圆 PDE 上建立完整的算法验证链条}。这一路径最符合本科毕业论文的工作量与可操作性。可选取典型变系数 Poisson 方程或扩散方程作为模型，先实现高精度 FEM 快照生成，再分别实现 batch SVD 构造的 POD/RB 作为基线，以及 weighted inner product 下的增量 core SVD 作为对照。随后系统比较三类指标：近似误差、正交性指标、计算时间与内存占用 \cite{hesthaven2015,fareed2018,zhang2022}。

第二，\textbf{围绕 weighted inner product 下的数值稳定性展开更细致研究}。有限元背景中的自然内积不是标准内积，因此单纯移植一般数据流领域的增量 SVD 算法并不充分。一个值得重点考察的问题是：不同 $W$-正交化策略对误差、耗时和鲁棒性的影响如何。例如，可比较完全重正交、按阈值触发重正交、仅对新增方向递归正交化等不同策略，并用
\[
E_W=\|I-Q^TWQ\|
\]
作为统一指标衡量正交性保持效果。 \cite{fareed2020,zhang2022}。

第三，\textbf{研究截断策略与基维数控制}。增量 SVD 的价值并不只在“更新快”，还在于能边生成、边截断、边控制秩增长。可重点分析两类阈值：一类是残差阈值 $p<\mathrm{tol}$，用来判断新快照是否基本落在当前子空间中；另一类是最小奇异值阈值 $\sigma_k<\mathrm{tol}$，用来决定是否丢弃弱模态。不同阈值会同时影响误差、基维数、在线效率和正交性，因此值得系统实验比较 \cite{brand2002,fareed2018,zhang2022}。

第四，\textbf{从算法层面探索比 Brand 原始更新更稳健的结构}。这部分未必要提出全新的理论算法，但至少可以在实现和比较层面回答一个非常具体的问题：Brand 直接更新、Brand 间接更新、Zhang 型改进更新，在 PDE 快照数据上谁更适合做 POD/RB 离线构基。由于这一方向同时连接了计算效率、正交性保持与加权内积下的数值稳定性，因此即便论文最终侧重数值比较，也具有明确的方法学意义 \cite{brand2006,zhang2022}。

第五，\textbf{将研究结果进一步连接到更复杂的 many-query 任务}。如果时间允许，可以在论文结尾讨论其向更复杂模型的推广，包括抛物型方程、时变 PDE、非线性降阶，以及更复杂的采样机制如贪婪算法和自适应参数选点。这样做的意义在于说明：增量 SVD 并不只是一个线性代数技巧，而是 many-query 降阶计算链条中的关键中间环节，它直接影响离线构基效率，并最终影响在线快速求解的可行性 \cite{hesthaven2015,patera_rozza,zhang2022}。

综上，基于增量奇异值分解的 PDE 模型降阶方法并不是简单地把 SVD 换成增量 SVD，而是在 many-query、RB、POD、weighted inner product、数值正交性与截断控制之间建立起一套完整的方法链。已有研究表明，增量 SVD 在离线构基效率方面具有明显潜力，但其数值稳定性，尤其是加权内积下的正交性保持问题，仍然是进一步研究的关键所在 \cite{brand2002,brand2006,fareed2018,zhang2022}。


% 参考文献
{
\newpage
\bibliographystyle{gbt7714-numerical}
\phantomsection
\subsection{参考文献}
{\normalfont\CJKfamily{simsun}\zihao{5}\setlength{\baselineskip}{14pt}
\renewcommand{\refname}{\vspace{-\baselineskip}}
\bibliography{refs/01_Review}}
}